\begin{abstract}
We evaluate a white-box provenance prototype that associates a classifier checkpoint with a signed release manifest. The manifest records a SHA-256 Merkle commitment to the ordered 50,000-image CIFAR-10 training set and is signed with ML-DSA-65. A keyed weight projection embeds the Reed--Solomon encoding of the 256-bit manifest digest. A trusted auditor uses channel assignment and positive-scale fitting as extraction retries, and reports residual codeword agreement when exact verification fails. Across five paired seeds on SmallCNN, mean clean test accuracy was 84.60\% for the unwatermarked baseline, 82.37\% for binary cross-entropy (BCE) with RS32, and 84.10\% for hinge loss with RS32, a gain of 1.73 percentage points over BCE/RS32. Robustness testing reused three trained models per design across 252 checkpoint evaluations per design. A same-layer 512-bit overwrite caused rejection within a five-point accuracy-drop budget on all three hinge models. Neither BCE configuration recorded a successful removal in the tested suite. The verifier returned exact matches on all 189 evaluations of the specified permutations and input scalings, including combined edits. Two RS32 control cohorts, with 724 extraction cases and 54 attacked baseline evaluations, yielded zero detections. These observations describe the tested checkpoints and settings.
\end{abstract}
